\documentclass[a4paper]{scrreport}

\usepackage[T1]{fontenc}
\usepackage[utf8]{inputenc}
\usepackage[nswissgerman]{babel}
\usepackage[babel,german=swiss]{csquotes}
\usepackage[onehalfspacing]{setspace}
\usepackage[style=verbose-ibid]{biblatex} % -ibid


% Literaturverzeichnis hinzufügen
\addbibresource{bibliography.bib}


% Einstellungen für Darstellung der Fussnoten
\usepackage[bottom,norule,hang,perpage]{footmisc}
\setlength{\footnotemargin}{2mm} % Spacing in Fussnoten zwischen der Zahl und dem Text
% \renewcommand{\footnoterule}{\rule{45mm}{0.1mm}{\vspace*{5mm}}} % Trennlinie über Footer


% Einstellunge für Zitierart
\renewcommand*{\labelnamepunct}{\addcolon\addspace} % Interpunktion in Zitaten zwischen Autor und Titel
\renewcommand{\newunitpunct}{\addcomma\addspace} % Zitatelemente mit einem Komma trennen (nicht der standardmässig Punkt)
\renewcommand{\subtitlepunct}{\addperiod\addspace} % Interpunktion in Zitaten zwischen Titel und Untertitel


% Spacing des Papiers
\usepackage{geometry}
\geometry{a4paper,top=30mm,right=27mm,bottom=30mm,left=27mm}


% Überschriftentiefe die im Inhaltsverzeichnis nummeriert werden
\setcounter{secnumdepth}{3}


% Überschriftentiefe die im Inhaltsverzeichnis angezeigt werden
\setcounter{tocdepth}{3}


% Abbildungen (figure) innererhalb einer Section durchnummerieren
% https://tex.stackexchange.com/a/425603
% \let\counterwithout\relax
% \let\counterwithin\relax
% \usepackage{chngcntr}
% \counterwithin{figure}{section}
% \counterwithin{table}{section}


% Sections durchgehend nummerieren und nicht von Chapters unterbrechen
% \usepackage{chngcntr} % chngcntr bereits eingebunden
% \counterwithout{section}{chapter}


% Package um Links im PDF-Dokument zu verwenden
\usepackage{url}


% Package um Bilder (PNG) und PDFs als figures einzubinden
\usepackage{graphicx}


% Placeholder for figures: \includegraphics[width=.5\textwidth]{example-image-a}
% https://tex.stackexchange.com/a/413186
\usepackage{mwe}


% Package um Symbole für Musik bereitzustellen: \flat \sharp
% \usepackage{textcomp}


% Package um Symbole Noten bereitzustellen: \quarternote \twonotes
% \usepackage{wasysym}


% Package um Generalbassbezifferungen im Fließtext besser zu verwenden
% https://tex.stackexchange.com/a/267644
% https://tex.stackexchange.com/a/267618
% \usepackage{amsmath} % \overset{above}{main}


% Command für inline Code im Fließtext
\newcommand{\code}{\texttt}


% Package und Command um Text einzukreisen für Bassstufen
% https://tex.stackexchange.com/a/199056
\usepackage{tikz}
\newcommand*\circled[1]{\tikz[baseline=(char.base)]{\node[anchor=text, shape=circle,draw, inner sep=0pt, minimum size=1.2em] (char) {#1\strut};}}


% Git Hash der Version erkennen und auf \gitAbbrevHash legen
\usepackage{xstring}
\usepackage{catchfile}
\CatchFileDef{\gitHash}{.git/refs/heads/master}{}
\newcommand{\gitAbbrevHash}{\StrLeft{\gitHash}{7}}


% Alle Quellen im Literaturverzeichnis anzeigen, auch nicht verwendete
\nocite{*}


% Einstellungen für das Stichwortverzeichnis
% \usepackage{makeidx}
% \makeindex
% \renewcommand{\indexname}{Stichwortverzeichnis}


% Nummerierung von Sections ohne die Zählung von Chapter
\renewcommand*\thesection{\arabic{section}}


% Spacings in Tabellen
\setlength{\tabcolsep}{1em} % horizontal padding
{\renewcommand{\arraystretch}{1.5}% vertical padding


% Fussnoten in Tabellen aktivieren mit dem Command \tablefootnote
\usepackage{tablefootnote}


% Zeilenumbrüche in Tabellen erlauben mit \thead{Mixo-\\lydisch} und |makecell{}
% https://tex.stackexchange.com/a/176780
\usepackage{makecell}


% Fügt Fraktur Schriftarten hinzu: \textfrak{Fraktur} \textswab{Schwabacher} \textgoth{Gothic}
% \usepackage{yfonts}


% "Abbildung ..." unter den figures umbenennen
\usepackage[figurename=Notenbeispiel]{caption}
\renewcaptionname{nswissgerman}{\listfigurename}{Verzeichnis der Notenbeispiele}


% Gestrichelte horizontale Linien in Tabellen mit \hdashline
\usepackage{arydshln}


% Abstände im Inhaltsverzeichnis zwischen Zahl und Titel anpassen
\usepackage{tocloft}
\setlength{\cftsecnumwidth}{1.5em}
\setlength{\cftsubsecnumwidth}{2.1em}
\setlength{\cftsubsubsecnumwidth}{2.75em}


% Sicherstellen, dass figures in der section ausgegeben werden zu der sie gehören
\usepackage[section]{placeins}




\title{Gemeinsam musizieren Lernen}
\subtitle{Master Musikpädagogik}
\subject{Masterthesis}
\titlehead{
	\hfill
	\includegraphics{img/Logo_Musikhochschule.pdf}
}
\author{
	Katharina Jäckle\\
	Hochschule für Musik Freiburg
}
\dedication{
	{
		\Huge
		Gemeinsam musizieren Lernen\\
	}
	\vspace{0.75cm}
	{
		\Large
		Das gemeinsame Musizieren Lernen\\
		oder auch gemeinsam das Musizieren Lernen\\
	}
	% \vspace{1.5cm}
	% {
	% 	\normalsize
	% 	Ausserdem noch eine kleine Danksagung \\
	% 	oder weitere kurze Unteruntertitel hier\\
	% }
	\vspace{12cm}
	{
		\footnotesize
		Text- und Notensatz mit \LaTeX, LilyPond und git.\\
		\small{Version: \gitAbbrevHash}
	}
}
\publishers{
	Gutachter: Andreas Doerne\\
}
\date{\today}



\begin{document}

	\maketitle

	\tableofcontents

    \addchap{Vorwort}

Danke an die Einflüsse, Anstöße und den Input meiner Professoren Wolfgang
Lessing und Andreas Doerne, die vorallem im musikpädagogischen Bereich prägend
waren aber sich auch auf mein künstlerisches Dasein ausgewirkt haben. 
An Andreas Doerne, der unsere Masterarbeit betreut
Wolfgang Lessing, der uns zum Wettbewerb gebracht hat etc.
Michael Stecher?
Und an meinen Kommilitonen Timo Langpap, mit dem ich Co-Working betrieben, mich
immer wieder ausgetauscht und diskutiert habe. Der außerdem mit seiner Biografie
noch einmal einen anderen Blick auf die Dinge hat und auch in seinem Genre dem
Jazz und Pop auch andere Gepflogenheiten gewohnt ist, was immer zu einem
spannenden Austausch geführt hat. 

Außerdem bedanke cih mich fürs Korrekturlesen und die Unterstützung in der
Nutzung des Programms LaTex bei Wolfgang Drescher, der mir bei all diesen Fragen
stets zur Seite stand. 

Durch unsere Gedanken zu einem eigenen Musikschulmodell, habe ich mir auch über
das gemeinsame Musizieren und Musiklernen in der Gemeinschaft viele Gedanken
gemacht.




\vspace{0.5cm}

\begin{flushright}
	{
		\small
		Freiburg im Breisgau, September 2021\\
		--- Katharina Jäckle
	}
\end{flushright}

	\addchap{Einleitung}
Das gemeinsame Musizieren zieht sich bei meinem Lebenslauf wie ein roter Faden
durch die ganze Laufbahn. Als Instrumentallehrerin habe ich allerdings auch
beobachtet, wie es für Schüler ist, die lediglich den
Einzel-Instruemntalunterricht kennen. Dabei ist mir aufgefallen, dass ein
wichtiger Bestandteil der für mich die Musik ausmacht fehlt. Nämlich das
gemeinsame Musizieren. Wenn man alleine den Aspekt des
Differenziellen Lernens betrachtet wird schnell klar, das Einzelunterricht
alleine nicht ausreicht um einen Schüler motiviert am Instrument zu halten und
ihn weiter zu bringen. Da sich auf der sozialen Ebene im Miteinander so viel
ereignet, kann eine einzelne Lehrkraft nicht alle Ebenen abdecken und bedienen.
Oft ist in unserem Musikschulsystem das gemeinsame Musizieren aber gar nicht
vorgesehen, sondern vielmehr ein optionales Zusatzangebot, im Orchester, im Chor
oder in der Band zu spielen. Dieses Zusatzangebot ist auch nicht immer gegeben,
je nach Größe der Musikschule und schon sind die Schüler schnell auf sich und
ihren Instrumentallehrer allein gestellt. In anderen Ländern gibt es andere
Sitten und somit auch verschiedene Systeme, in denen Instrumentalschüler mehr im
Miteinander das Instrument lernen. Einige will ich deshalb näher beleuchten und
unteruchsen, inwieweit sie für die Lernenden wertvoll sind. Im Vordergrund
meiner Arbeit soll aber stehen, warum das gemeinsame Musizieren so
wichtig ist und das werde ich ich im Folgenden beleuchten. 


"(...) dass Muisk in ihrem Wesen Mitteilung, Kommunikation, Gestaltung, das
Vermitteln einer Botschaft ist, die in Sternstunden über das Erleben der
Einzelnen hinausreicht und sie miteinander verbinden kann." S. 11

"Musik als Kommunikationsmittel" S. 17

Aristoteles: "der Schlüssel zu einem glücklichen Leben besteht im Miteinander,
das Leben in der Gemeinschaft unter Freunden." (S. 9 "Die Kunst zu unterrichten)

"Sie muss gelernt und ein Leben lang geübt werden." (auch S. 9)


"Musik ist Teil der Symbolwelt des Menschen: Mitteilung, Kommunikation,
Interaktion." Wolfgang Suppan, (S. 91 Andreas Doerne)


	\addchap{Gemeinsames Musizieren}

Das gemeinsame Musizieren zu lernen ist ein wichtiger Prozess, der in der
Instrumentalausbildung ein fester Bestandteil sein sollte. Schließlich müssen
die Instrumentalschüler*innen erst einmal an die Gruppenprozesse herangeführt
werden um zu lernen wie man (musikalisch) interagieren kann. Das schließt
allerdings nicht aus, dass dabei das gemeinsam das Musizieren an sich gelernt wird. Im
besten Fall wird nämlich genau das erreicht. Bei dem gemeinsamen Musizieren
steht die Freude am Musizieren im
Vordergrund und nicht die technischen Hürden, die ein*e
Instrumentalschüler*in überwinden muss. Insofern sollten am
besten der Einzelunterricht und das gemeinsame Musizieren in der Praxis ineinander greifen.

Bei dem gemeinsamen Musizieren spielen viele verschiedene Faktoren eine Rolle.
Im Folgenden werde ich deswegen mehrere Aspekte beleuchten und erörtern.
Beispielsweise welche Aspekte die Schüler*innen weiterbringen und für sie
gewinnbringend sein können. Außerdem werfe ich einen Blick auf die Aufgabe der
Instrumentallehrer*innen um darauf einzugehen, wie man verschiedene Gruppen und
Situationen wirkungsvoll leiten kann. Alle Aspekte spielen sich auf der sozialen
interaktiven Ebene ab, im Wechselspiel zwischen dem Individuum und der Gruppe.
Dafür werde ich die Gruppenprozesse, sowie mögliche Aktionsformen für Gruppen,
die Hierarchie innerhalb einer Gruppe, die Kommunikation, Motivation und das
psychosoziale Lernen, welches in der Gruppe stattfindet betrachten, sowie die
Bedeutung von Gruppenarbeit für den Lernenden und die Bedeutung für
die Lehrkraft.

















%"Musik als Kommunikationsmittel" S. 17 doerne nicht

%Aristoteles: "der Schlüssel zu einem glücklichen Leben besteht im Miteinander,
%das Leben in der Gemeinschaft unter Freunden." (S. 9 "Die Kunst zu
%unterrichten)

%%"Sie muss gelernt und ein Leben lang geübt werden." Zitat? doerne nicht


%"Jedes Menschen Persönlichkeit, das bedeutet: seine Fähigkeiten, seine Art zu
%denken und zu fühlen, wird durch Umstände und Umgebung geschnitzt und
%gemeißelt." (S. 20 shinichi Suzuki)

%Andreas Doerne schreibt sehr zutreffend "(...) dass Musik in ihrem Wesen
%Mitteilung, Kommunikation, Gestaltung, das Vermitteln einer Botschaft ist, die
%in Sternstunden über das Erleben der Einzelnen hinausreicht und sie miteinander
%verbinden kann." (Nope niht doerne) zitat?



\section{Gruppenprozesse}


- Individualität in der Gruppe 
"(...)Individualität zeigt sich gerade auch im Umgang mit anderen,
Individualität entwickelt sich im Vergleich zu und in ABgrenzung von anderen
oder auch dadurch, dass Positionen in einer Gruppe besetzt werden."

- unterschiedliche Alters- und Lernniveaus (manchmal) evtl. S. 95

- Lernen / Differenzielles Lernen? 
"Ein Großteil musikalischen Lernens findet über Imitation statt (...) S. 98

- Hierarchie
fortgeschrittenere Schüler zu Mentoren erklären und sie damit mitverantwortlich
für deren musikalische Entwicklung und den Fortschritt des Ensembles zu machen.
S. 95

- Altersgruppen/Erwachsene
"Wiedererlangen sozialer Integration scheint ein wesentliches Motiv
musikalischer Betätigung bei Erwachsenen (...). " s. 118

- "Selbsttätigkeit" (die Kunst zu unterrichten)
"Erst wenn es eine Notwendigkeit für einen Schüler gebe, eine Arbeit wirklich
selbst zu vollbringen, würde er ganzheitlich mit seinen Gefühlen, Trieben und
Willenskraft in die Tätigkeit involviert. (S. 63 bzw Eigl von Kerschensteiner)

-Gruppe als Korrektiv

\subsection*{Bedeutung für den Schüler*in}

Welche Unterrichtsform für welchen Schüler besonders geeignet ist, ist ein sehr
individuelles Unterfangen. Während der eine Schüler die ganze Aufmerksamkeit
eines Lehrers genießt, kann sich der andere im Einzelunterricht unwohl fühlen.
Wer ein ser kommunikativer Mensch ist, ist hingegen vermutlich in den
Gruppenunterrichten besser aufgehoben. Oftmals ändern sich die Unterrichtsformen
für einen Schüler allerdings auch, da es gängig ist in Gruppenunterrichten zu
starten. Die Form eines Kleingruppenunterrichts bei Anfängern bietet sich auch
an, siehe Zitat \autocite[220]{Busch:grundwissen_Instrumentalpädagogik}

"Man erfährt beim Musizieren soziale Resonanz(...)" S.28

"Eine Steigerung musikalischer Kommunikation erfährt der Schüler im
Zusammenspiel mit dem Lehrer, im Gruppen- oder Ensembleunterricht. (...) Die Individualität des eigenen Musizierens entwickelt sich in
Beziehung zu
anderen."  S.99 Punkt 5

"Das Gefühl von Zugehörigkeit und Zusammenklang beim Spielen im Ensemble."


\subsection*{Bedeutung für den Lehrer}

Je nach Gruppengröße muss der Lehrer verschiedene Aufgaben übernehmen, zu der
auch eine unterschiedliche Vorbereitungsnotwendigkeit hinzukommt. Während der
der Lehrende im Einzelunterricht unmittelbar auf die Impulse des Schülers
eingehen kann, so muss der Lehrende bei Kleingruppen neben klaren Vorgaben die
Schüler zu einem von- und miteinander Lernen anregen und sich selbst auch einmal
zurücknehmen. Im Großgruppenunterricht ist es essentiell, dass der Lehrer im
Vorfeld eine Struktur und Stücke festlegt um die Inhalte gut anleiten zu können.\autocite[220]{Busch:grundwissen_Instrumentalpädagogik}
Meistens hat man in dieser Konstellation die Situation eines Frontalunterrichts,
was beim Klassenmusizieren von den Instrumentalpädagogen verlangt wird, während die
Instrumentalpädagogen mit diesem Szenario oftmals nicht vertraut sind.

"Für den Instrumentallehrer stellt Gruppen- oder Ensembleunterricht eine höhere
Herausforderung dar, da er seine Aufmerksamkeit (gerecht) und flexibel
verteilen, differenziertere Spielanweisungen geben und eine musikalisch und
menschlich komplexer Situtation steuern muss. Andereseits ergeben sich im Umgang
mit der Gruppendynamik auch andere Möglichkeiten der Motivation und Förderung
musikalischen Lernens." S. 99 Punkt 6

- (differenzielles Lernen)
"Von einem bestimmen Alter an ist den Heranwachsenden die Rückmeldung der
Gleichaltrigen sowieso wichtiger als die des Lehrers, dann bewirtk die Kritik
eines Pultnachbarn manchmal eher eine Veränderung oder eine intensivere
Auseinandersetzung mit einem Stück als die des Dirigenten "Da vorn." S. 95


Evtl. ins Fazit? ((Die Musikhochschule müssten in diesem Punkt noch etwas besser ausbilden, da
jeder Instrumentalpädagoge darauf vorbereitet sein sollte und es so wichtig ist.
S. 94 oberer Abschnitt ))

"(...) oft ziemlich eng ausgebildeten Instrumentalpädagogen(...) S. 119






	\addchap{Analyse verschiedener Systeme}


\section{Klassenmusizieren}





	\addchap{Der gängige Weg}

Wie in Deutschland der gängige Werdegang eines Instrumentalschülers aussieht,
werde ich im Folgenden beschreiben. Viele Schüler*innen kommen in ihrem Leben
ausschließlich in den Genuss des Einzel-Instrumentalunterrichts, aber es gibt in
der musikalischen Früherziehung, sowie auch ergänzend zum Einzelunterricht auch
in hierzulande einige Angebote, die Schüler*innen über den Einzelunterricht
hinaus wahrnehmen können und auch teilweise nutzen.


\section{Musikalische Früherziehung}

Das Arbeiten mit Gruppen ist ein grundlegendes Merkmal der Elementaren
Musikpädagogik (EMP). Wir können aus der EMP auch für andere Musiziergruppen
einige wichtige Bausteine ableiten, zumal sich das Elementare Musizieren an
Menschen jeden Alters wendet.
\autocite[226]{busch:grundwissen_instrumentalpaedagogik} Oft wird sie jedoch
besonders für die musikalische Früherziehung genutzt und richtet sich bei vielen
Angeboten an junge Familien oder Kindergartenkinder. Es ist sehr wichtig, dass
die jüngste Generation aus unserer Gesellschaft schon im frühen Alter die
Gruppenprozesse und die Musik kennenlernt, da sie auch in ihren sehr jungen
Jahren schon sehr aufnahmefähig sind. Wieviel sie sich in den ersten
Lebensjahren aneignen zeigt sich alleine darin, dass sie ungefähr im Alter von
vier Jahren der Muttersprache mächtig sind. "Auf die Musik übertragen sind sich
alle großen Musikpädagogen von Zoltan Kodaly bis Shinichi Suzuki einig, dass ein
Kind unbedingt so früh wie möglich mit der Musik in Kontakt kommen muss. "In
diesem Sinne gilt es, ein Kontinuum der der, musikalischen Muttersprache zu
entwickeln." \autocite[45]{ernst:die_zukunftsfaehige_musikschule} Nicht nur für
die Gruppenprozesse, sondern vor allem auch für das natürliche heranführen an
die Musik, ist die Musikalische Früherziehung besonders wichtig. Der Ort dieser
musikalischen Basisbildung sollte nicht erst in der Musikschule sein, sondern
bereits im Kindergarten oder vorher.
\autocite[43]{ernst:die_zukunftsfaehige_musikschule}Der Vorteil ist, dass das
Das Musizieren im Kindergarten ist zum Einen für jedes Kind zugänglich und
findet zum Anderen auch im Alltag integriert statt. Im Anschluss können die
Schulkinder außerdem durch eine bessere Orientierung an den Instrumenten, das
Instrument auswählen, das ihnen am meisten zusagt. Diese Orientierungsphase ist
nicht zu unterschätzen und besonders heutzutage wichtig, wo die Kinder nicht
mehr vielfältig musizierend, sondern einförmig konsumierend mit der Musik
aufwachsen.\autocite[37]{ernst:die_zukunftsfaehige_musikschule}

Neben dem Umgang mit Gruppen werden im Folgenden die weiteren Grundgedanken der
EMP dargelegt. Ganz vorne weg ist die wichtige entwicklungspsychologische
Erkenntnis nach Oerter und Lehmann zu betonen, die besagt, dass jeder Mensch
musikalisch ist.\autocite[88]{musikalische_begabung} Geht man von diesem
Grundgedanken aus, so soll die EMP ermöglichen, dieses musikalische Potenzial zu
entwickeln, wobei die Entfaltung der musikalischen Ausdrucksfähigkeit die
Persönlichkeitsbildung des Menschen unterstützt. Das gelingt durch das
musikalische Gestalten, wozu verschiedene Handlungsbereiche zählen, wie das
sensorische Sensibilisieren, die Bewegung, das Singen, Sprechen, das Lernen die
Zeit zu strukturieren, (Elementar-  und Körper-)Instrumente spielen, das
Visualisieren sowie das Hören.
\autocite[227]{busch:grundwissen_instrumentalpaedagogik} Als grundlegendens
Musikverständnis bilden in der EMP Musik und Bewegung stets eine Einheit.

Zwei übergeordnete Erziehungsziele sind es, die psychologischen und sozialen
Ziele zu schulen. Was die beiden aber verbindet und vereint ist die Entwicklung
von Selbstwahrnehmung und sozialer Beziehungsfähigkeit, durch die zunehmend die
Ich-Identität ausgebildet wird. Zu den psychologischen Zielen gehört es,
Geborgenheit zu erfahren und die Selbstwahrnehmung zu fördern. Außerdem fördert
die EMP die "Ich-Erweiterung" durch Identifikationsprozesse mit verschiedenen
Lebewesen, Naturerscheinungen und unbelebten Objekten. Es geht des weiteren bei
den psychologischen Zielen darum Ideen zu entwickeln und persönlich
herausgefordert zu werden, in dem sich die Schüler präsentieren müssen und
dadurch Unsicherheiten entstehen können. Es gilt zu Lernen Spannungen
auszuhalten, das Durchhaltevermögen zu steigern und Befriedigungsaufschub
ertragen zu lernen. Ein weiterer Aspekt ist es Initiative zu ergreifen und
Entscheidungen zu fällen, Beziehungen zur Gruppe oder einem Partner aufzubauen
und als Letzten Aspekt auch das Zurücktreten und Verzichten, sowie
Selnbstkontrolle zu lernen. Dies sind alles psychologische Effekte, die im
Umgang mit der Gruppe besonders geschult und gelernt werden können; besser als
im Einzel- Instrumentalunterricht. Nichts destotrotz beziehen sie sich auf das
Individuum selbst.

Die sozialen Aspekte hingegen beziehen sich, wie schon erwähnt auf die Gruppe.
Hier geht es darum, die Anderen wahrnehmen zu lernen. Dazu zählt auch
aufeinander zu reagieren, in Kommunikation zueinander zu treten und aufeinander
Rücksicht zu nehmen. Das Abwarten und eine Reihenfolge einhalten, wird auch
geschult, indem sich die Gruppe selbst regeln lernt und ein Gruppengefühl
entwickelt wird. Das funktioniert wiederum nur, wenn die Schüler anfangen sich
mitverantwortlich zu fühlen und auch Ideen andere aktzeptieren.

Um sich diesen Zielen zu nähern, geht die EMP spielerich mit Individual und
Kollektivgestaltung um und versucht durch vielfältige Aktionsformen verschiedene
Interaktionen zu ermöglichen. Dadurch geschieht die soziale Sensibilisierung, zu
der noch die auditive Sensibilisierung hinzukommt. 



\subsection{Aktionsformen in der Gruppe und Persönlichkeitsentwicklung}
Die Sozialform der EMP ist wie eingangs erwähnt immer die Gruppe, die aber
wiederum verschiedene Spielräume für verschiedene Aktionsformen ermöglicht. So
kann es Phasen geben, in denen jeder für sich tätig ist, Partnerarbeit oder die
Arbeit in der Gruppe als Ganzes. Alle Formen basieren auf dem
beziehungsorientierten Arbeiten, in dem "über die Beziehungen, die zur Musik
aufgebaut werden sollen, [...] auch die Kontakte mit anderen Gruppenmitgliedern
im Fokus." \autocite[10]{dartsch:kern_des_musizierens} stehen. Mit dem
ganzheitlichen Menschen im Blick soll neben der angestrebten
Persönlichkeitsbildung die soziale Kommunikation ausgebaut werden. "Dass viele
dieser Kinder gut mit Musik kommunizieren können und damit sowohl im oft
notwendigen Gruppenunterricht, aber auch bei musikalischen Dialogen mit dem
Instrumentallehrer ausdrucksfähiger sind." S. 25%Andreas oder erstes Buch? s. 25





%"(...) sind viele der damit verbundenen Techniken und Methoden für den
%"Ernsthaften" Instrumentalunterricht durchaus hilfreich." S. 25 Anselm Ernst S.
%37 ff




%"Für die Musikschule reklamieren wir den Erziehungs- und Bildungsauftrag, in
%allen Schichten der Gesellschaft so umfassend wie möglich eine Musizierkultur
%aufzubauen" (S. 37/38) Anselm Ernst




\section{Instrumentalunterricht}

Jeder Mensch, der ein Instrument erlernt, erhält in seinem Leben früher oder
später Instrumentalunterricht, was in Deutschland erst einmal Einzelunterricht
bedeutet. Der Instrumentalunterricht bietet aber viele verschiedene Formen vom
Einzelunterricht (EU), über den Gruppenunterricht (GU), bis hin zu dem
sogenannten \emph{Multidimensionalen Unterricht}. Man könnte auch soweit gehen,
dass Band- Chor- oder Orchesterproben mit einem Lehrer zu dem normalen
Instrumentalunterricht zählen. Als Erstes verknüpfen wir im europäischen System
jedoch mit dem klassischen Instrumentalunterricht den Einzelunterricht, wo eine
Lehrkraft eine/einen einzelnen Schüler*in im Normalfall zwischen dreißig und neunzig
Minuten unterrichtet. Dazu gehe ich im Folgenden ein. Was den Gruppenunterricht
ausmacht, wird auch näher beleuchtet werden. Fest steht auf jeden Fall, dass im
Idealfall das Lehren und Lernen flexibel gestaltet werden und sich im besten
Fall der EU und GU ergänzen. In unserem Musikschulsystem ist leider noch zu
wenig Raum für die Kombination aus individuellem- und gemeinschaftlichen Lernen.
So schreibt Anselm Ernst in seinem Buch Die zukunftsfähige Musikschule:
\enquote{Wer genügend Raum und Zeit hat, kann sich mit Ruhe, Geduld und
Gelassenheit den Menschen und den Sachen widmen.}
\autocite[84]{ernst:die_zukunftsfaehige_musikschule} Bei allen
Unterrichtsformen, so unterschiedlich sie sein mögen und mit ihren jeweiligen
Vor- und Nachteilen soll es darum gehen, dass der/die Instrumentalschüler*in das
für ihn am besten geeignete Lernumfeld geboten bekommt. Dazu sollte es auf alle
Fälle variabel sein und keine starre einseitige Unterrichtsform, sondern eine
Wechselwirkung aus allen Möglichkeiten sein, aus denen wir Musiker schöpfen
können. 


\subsection{Einzelunterricht}

Im Einzelunterricht (EU) erhält der/die Instrumentalschüler*in die volle
Aufmerksamkeit der Lehrkraft. Dies hat seinen Vorteil darin, dass die Lehrkraft
auf der persönlicheren Ebene viel mehr auf den/die Schüler*in eingehen kann,
sowie spieltechnisch in einem auf den/die Schüler*in abgestimmten Tempo arbeiten
kann. Außerdem erlaubt der EU, tiefer und detaillierter auf das ein oder andere
Thema einzugehen, zum Beispiel in der Spieltechnik oder Körperhaltung. Die volle
Aufmerksamkeit der Lehrperson kann jedoch auch genauso gut zum Nachteil werden,
da der/die Schüler*in ausschließlich von der Lehrperson seine Anweisungen erhält
und der Lernweg ausschließlich eindimensional stattfindet. Der/die Lernende kann
sich zwar etwas von der Lehrperson abschauen und etwas Unausgeprochenes aus dem
Unterricht mitnehmen, aber die Lehrperson alleine kann gar nicht soviel bieten,
wie eine Gruppe von Mitschüler*innen. Gerade für Schüler*innen, die ein Problem
mit der Hierarchie haben, ist möglicherweise der Einzelunterricht nicht das
optimale Lernumfeld. Andere Schüler hingegen genießen es, endlich einmal die
volle Aufmerksamkeit einer Erwachsenen Person und gegebenenfalls sogar
Vertrauensperson zu bekommen. Diese Erfahrung habe ich besonders bei meinen
Internatsschüler*innen gemacht, die zu Hause viele Geschwister haben. 

\subsection{Gruppenunterricht} 
In den Musikschulen in Deutschland wird für Anfängergruppen oft ein
Gruppenunterricht (GU) angeboten, aber früher oder später haben die
Schüler*innen Einzelunterricht, oder aufgehört, das Instrument zu spielen. Beim
GU muss man zwischen Klein- und Großgruppen unterscheiden. In einem
Kleingruppenunterricht sprechen wir von zwei bis fünf Schülern, während ein
Großgruppenunterricht alles darüber hinaus beschreibt.
\autocite[219]{busch:grundwissen_instrumentalpaedagogik} Laut Anselm Ernst gilt
ein Gruppenunterricht ab drei oder mehr Schülern, da man erst ab drei Schülern
von einer Gruppe sprechen kann.
\autocite[79]{ernst:die_zukunftsfaehige_musikschule}
Allerdings gibt es zusätzlich noch das Szenario von zwei Schülern, die gemeinsam
Unterricht haben und somit weder dem Einzel-, Klein- noch dem Gruppenunterricht
zugeordnet werden können. Hier sprechen wir dann vom
Partnerunterricht.\autocite[219]{busch:grundwissen_instrumentalpaedagogik} Bei
zwei Schüler*innen dominiert der Lehrer trotz allem noch das Geschehen, aber
nichts desto trotz können die Schüler gemeinsam Duos spielen, ohne, dass die
Lehrperson direkt als Duo-Partner involviert ist, wie im EU. \enquote{Während im
Großgruppenunterricht und somit auch im Klassenmusizieren bzw. Klassenunterricht
dem Erwerb spieltechnisch-musikalischer Fertigkeiten Priorität zukommt, ist das
Spielen im Ensemble hiervon anzugrenzen, da dies das gemeinschaftliche
Musizieren fokussiert und in der Regel einen vorausgehenden oder parallel
stattfindenden Instrumentalunterricht
voraussetzt.}\autocite[219]{busch:grundwissen_instrumentalpaedagogik}

Die Vorteile des Gruppenunterrichts sind sehr vielseitig. Abgesehen davon, dass
gemeinsames Musizieren einfach mehr Spaß macht und da kann ich aus Erfahrung von
meinen Schüler*innen sprechen, die förmlich danach lechzen, wieder gemeinsam zu
musizieren. Ein weiterer Aspekt ist, dadurch, dass das soziale Lernen in den
Vordergrund rückt, verschmelzen das soziale und musikalische Lernen. Außerdem
enstehen bei den Schüler*innen durch das soziale Miteinander, eine stärkere
Bindung zur Musik und darüber hinaus auch Freunschaften. Auch die Übemotivation
ist höher, da sich die Schüler*innen vor den anderen zum Einen nicht blamieren
wollen und zum Anderen auch eine Verantwortung für das Mittragen des Unterrichts
übernehmen. Außerdem können die Schüler*innen die Stücke gemeinsam Üben, sodass
sie auch beim Üben nicht alleingelassen sind und das Üben selbst auch gemeinsam
lernen. In einem Alter, wo das Instrumentalspiel vermeintlich uncool werden
könnte, lernen die Schüler*innen aber in der Gemeinschaft, dass es auch andere
Personen gibt, die ein Instrument spielen, wodurch sich die Schüler in ihrer
musikalischen Betätigung gegenseitig bestärken.
\autocite{ernst:die_zukunftsfaehige_musikschule}

Im Vergleich zum Musikunterrich in der Schule oder auch dem sogenannten
Klassenmusizieren, muss man außerdem den Vorteil sehen, dass die
Instrumentallehrer*innen sehr flexibel sind in ihrer Gestaltung, da sie sich an
keine Lehrpläne halten müssen. Dies gilt besonders, wenn das Musizieren im
Rahmen von der Musikschule stattfindet.



\subsection{Klassenmusizieren}
Eine Form des Gruppenunterrichtes in Großgruppen stellt das Klassenmusizieren
dar. Es findet oft in sogenannten Bläser- oder Streicherklassen statt, sodass
die Schüler*innen je nach Angebot der Schule auf die Familie der Streich- oder
eben Blasinstrumente festgelegt werden. Das Klassenmusizieren ist allderings
nicht zu verwechseln mit dem normalen Musikunterricht, welcher trotzdem noch
abgehalten wird. Im Idealfall läuft es so ab, dass die Schüler*innen die
verschiedenen Instrumente, die zur Auswahl stehen, kennen und anschließend eines
auswählen dürfen. Ziel ist es auch hier wieder, dass die Kinder im gemeinsamen
Musizieren, zu dem auch das Singen, Solmisieren und Rhythmusspiele gehören,
einen ganzheitlichen, praktischen Weg zur Musik finden.
\autocite[91]{ernst:die_zukunftsfaehige_musikschule} Dadurch, dass das
Klassenmusizieren im normalen Unterrichtsalltag integriert wird, lernen die
Schüler*innen das Instrument, wie jedes andere Fach, das sie unterrichtet
bekommen als selbstverständlich kennen. Dadurch, dass es im Stundenplan
integriert ist, ist es kein besonderes außerschulisches Angebot, welches nur von
wohlhabenden Eltern ermöglicht werden kann, sondern für Jeden zugänglich. So
haben die Schüler*innen die Gelegenheit, die Musik wie eine weitere Sprache
dazuzulernen. Im Unterschied zum beispielsweise Englisch-Unterricht, sind sie
aber nicht auf sich alleine gestellt, vor allem dann, wenn sie bei Prüfungen
ihre eigene Leistung unter Beweis stellen müssen. \enquote{Beim instrumentalen
Klassenunterricht jedoch verwandelt sich die Klasse der Schüler/innen in ein
Ensemble. Sie lernen gemeinsam und musizieren gemeinsam.}
\autocite[92]{ernst:die_zukunftsfaehige_musikschule} Wie Anselm Ernst außerdem
sehr zutreffend beschreibt, verbinden sich die Einzelleistungen und addieren
sich nicht. Dass die Schüler sich als Gemeinschaft erleben ist der zentrale Wert
dieses Musizierunterrichts!

%"Es wächst die Zahl der Kinder, die über Formen des Klassenmusizierens an
%%Grundschulen (...)) gewonnen werden." %/S. 118 welches buch? evtl.
%Doerne?-nope/ anselm ernst?

%Das Klassenmusizieren lehnt sehr stakr an das Model von JeKits an, oder auch
%umgekhert. 


%JeKids: meine Annahme, dass es Quatsch ist, einem Schüler im Klassenmusizieren
%ein Isntrument beizubringen.




\subsection{Multidimensionaler Unterricht} 
Das Konzept des Multidimensionalen Unterrichts weitet die Unterrichtsformen in
alle denkbaren Strukturen aus und ist derzeit noch nicht so weit in Deutschland
verbreitet. Das Konzept stammt von Gerhard Wolters und beschreibt eine gute
Möglichkeit, wie wir in unseren Musikschulen den Multidimensionalen Unterricht
als eine wunderbare Ergänzung zum Einzelunterricht etablieren könnten. Er
bezieht sieben Dimensionen ein, die sehr unterschiedliche Richtungen beinhaltet.
\autocite[86ff]{ernst:die_zukunftsfaehige_musikschule} Wolters spricht von der
Partnerschaft zwischen zwei und mehr Personen. Hier steht das Voneinander Lernen
im Vordergrung. Die zweite Dimension bezieht sich auf den Unterricht in mehreren
Räumen. Das bedeutet aber auch, dass die Lehrperson nicht immer anwesend sein
muss, während eines Unterrichtes. Egal ob es sich um eine Gruppe oder einen
Einzelunterricht handelt. Sehr revolutionär ist auch der Gedanke der flexiblen
Unterrichtszeiten. \enquote{Flexible (längere oder kürzere) Unterrichtszeiten
passen sich an die individuellen Lernbereitschaften, Lernrhythmen, Lerntempi
usw. an.} \autocite[87]{ernst:die_zukunftsfaehige_musikschule} Außerdem hat es
den positiven Nebeneffekt, dass der/die Instrumentalschüler*in am Ende sogar
längere Unterrichtszeiten hat, als in dem gewohnten EU. Eine neue Dimension ist
der Unterricht mit mehreren Lerhkräften. Hier können die Interaktionen der
Lehrkräfte untereinander angeregt werden, was das Unterrichten für die
Instrumentallehrer*innen durch die gegenseitigen Anregungen interessanter
gestaltet. Auch die Schüler*innen profitieren davon, da jede Lehrkraft ihre
eigenen Stärken mitbringt und die Schüler*innen direkt von verschiedenen Stärken
profitieren können. Was es in Kinderorchestern bereits gibt, wird in
Gruppenunterrichten, wenn sie denn angeboten werden, sehr selten praktiziert,
ist das Lernen mit Partnern unterschiedlichen Alters. \enquote{Intelektuell,
körperlich, emotional und sozial sind gleichaltrige Kinder/Jugendliche sehr
unterschiedlich entwickelt, so dass auch in dieser Dimension Flexibilität
erforderlich ist.} \autocite[87]{ernst:die_zukunftsfaehige_musikschule} Die
sechste Dimension schließt sich dem unterschiedlichen Alter an, denn Wolters
schreibt, dass auch Unterricht von Schüler*innen mit unterschiedlichen
Lernniveaus stattfinden soll. Zum Einen entspricht es der Lebensrealtität und
zum Anderen kann man die Schüler*innen in diesem Zusammenhang wunderbar dazu
anleiten, sich gegenseitig zu helfen, ihnen beibringen füreinander ein Vorbild
zu sein, sich gegenseitig zu motivieren und auch Verantwortung zu übernehmen.
\enquote{Von Gleichaltrigen oder Älteren lernen Kinder/Jugendliche oft besser
als von der Lehrperson.} \autocite[87]{ernst:die_zukunftsfaehige_musikschule}
Die letzte Dimension beschreibt den Unterricht mit verschiedenen Instrumenten.
Dadurch
\enquote{öffnet sich der Gruppenunterricht zum Ensembleunterricht.}
\autocite[87]{ernst:die_zukunftsfaehige_musikschule}


\section{Ensembles}

Nach dem Einzelunterricht, manchmal erst im späteren Leben und manchmal schon in
der frühen Kindheit, besuchen Instrumentalist*innen verschiedene Ensembles in
Ergänzung zum Instrumentalunterricht, oder auch unabhängig davon. Es gibt einige
Punkte, die alle Ensembles gemeinsam haben. So beudeutet das Wort
\emph{Ensemble} zunächst einmal \enquote{eine Gruppe von Dingen, die eng
zusammengehören.}\autocite{wikipedia:gruppe} Wenn man die Bedeutung weiter auf
die Kunst beziehen möchte, so steht in der Definition: \enquote{eine Gruppe von
Künstlern -wie Schauspieler, oder Musiker, die gemeinsam etwas vortragen.}
\autocite{wikipedia:gruppe} Das interessante an dieser Beschreibung ist, dass
das Wort \emph{vortragen} impliziert, dass diese Gruppe auch auf der Bühne oder
in einem anderen öffentlichen Rahmen mit ihrem Ensemble auftritt und das
Erarbeitete präsentiert. Betrachtet man den Einzel- Instrumentalunterricht, so
ist es nicht unbedingt selbstverständlich, dass der/die Schüler*in das
Erarbeitete öffentlich vorträgt. Gerade bei Privatschüler*innen werden oft keine
Vorspiele durchgeführt, da es ohne eine Institution im Hintergrund, die darauf
Wert legt, nicht immer dazu kommt, zumal auch nicht immer die räumlichen
Möglichkeiten dafür vorhanden sind. Wie sieht es jedoch mit dem
Gruppenunterricht aus? Er findet meistens im Zusammenhang mit einer
Institution statt, sei es dem Kindergarten, der Schule oder einer Musikschule.
Interessanterweise ist es dort, wenn auch teilweise unausgepsrochen,
selbstverständlich, dass die Gruppen bei einem besonderen Anlass oder auch bei
einem dafür geschaffenen Vorspiel, ihre Stücke darbietet. Andreas Doerne
beschreibt in seinem Buch \emph{Umfassend Musizieren} das Ensemble als
\enquote{eine musikalische Ganzeheit, die sich durch ein gleichberechtigtes,
dialogisches Miteinander und Zusammenspiel ihrer einzelnen Teile (Mitglieder)
konstituiert.} \autocite[62]{doerne:umfassend_musizieren} Diese Definition
trifft aber meiner Meinung nach nicht in allen Punkten auf ein Ensemble zu. So
wird im weiteren Verlauf der vorliegenden Arbeit deutlich, dass beispielsweise in einem Orchester der/die
Dirigent*in die Spielimpulse gibt und auch vorgibt, wie zu spielen ist,
weshalb in diesem Moment kein gleichberechtigtes Miteinander mehr existiert. Auch
das dialogische Miteinander ist nicht in jedem Ensemble erwünscht und
selbstverständlich. Wenn wir uns beispielweise die Suzuki-Talenterziehung
anschauen, werden wir sehr schnell feststellen, dass die Lehrkraft ganz deutlich
in der übergeordneten dominierenden Position ist, in der das dialogische
Miteinander auf einer sehr untergeordneten Ebene nur selten erwünscht ist.
Nichts desto trotz ist das \enquote{Ensemblespiel [ist] Kommunikation mit
musikalischen Körpern und Klängen.} \autocite[62]{doerne:umfassend_musizieren}
Sich dies zu vergegenwärtigen ist sehr hilfreich, da man schnell feststellt
welche zentrale Rolle dabei die Kommunikation spielt. Die Kommunikation wiederum
ist sehr eng mit der Psyche eines jeden verknüpft, wenn er Sitautionen für sich
selbst auswertet und mit anderen abgleicht und in vielerlei anderer Hinsicht.
Die soziale Komponente ist meiner Meinung nach das wesentliche Grundelement
eines Ensembles. Welche sozialen Eigenschaften innerhalb eines Ensembles eine
Rolle spielen ist im folgenden Zitat sehr zutreffend geschildert: \enquote{Das
Ensemble ist durchdrungen vom gegenseitigen Zuhören und Mitteilen, Sehen und
Gesehenwerden, Agieren und Reagieren, Führen und Geführtwerden, Hervorbringen
und Abnehmen von Spielimpulsen, insgesamt also vom wechselseitigen Geben und
Nehmen.} \autocite[62]{doerne:umfassend_musizieren} Diese sozialen Komponenten
finden wir auch in allen folgenden Formationen wieder. Außerhalb des Musizierens
treffen sich alle Ensembles zu mehr oder weniger regelmäßigen Proben und
veranstalten auch Probenwochenenden, an denen alle über die Musik hinaus
Zusammenkommen. Bei größeren Formationen gibt es oft auch außerhalb der Musik
einen Vorstand, der für die organisatorischen Fragen verantwortlich ist und bei
Problemen versucht die Gruppe zusammen zu halten. Das ist fast vergleichbar mit
unserer Gesellschaft, in der sich das Ganze nur in anderen Dimensionen abspielt.



\subsection{Orchester} 

Ein Sinfonieorchester ist, sofern man es überhaupt noch zu einem Ensemble zählen
kann, womöglich die größte Gruppierung von Musizierenden die zusammenspielen.
Aufgrund der Größe, trägt eine klare Hierarchie zur kleineren Unterteilung
innerhalb dieses großen Organismus bei. Auch aufgrund der Größe ist ein/eine
Dirigent*in, bei dem/der vorne der Klang zusammenläuft eine wichtige
Schlüsselfigur des Orchesters, da er/sie die Musik zusammenführen und regulieren
kann. Die einzelnen Spieler*innen sind mit ihrer eigenen Stimme beschäftigt und
sitzen teilweise sehr weit von anderen Stimmen entfernt, sodass die
Verantwortung bei dem/der Dirigent*in liegt, dies zusammenzubringen. Sowohl im
Timing als auch in der Dynamik. Das hat zur Folge, dass die Aufmerksamkeit der
einzelnen Spieler*innen neben den Stimmführer*innen auf den/die Dirigent*in
gerichtet ist, auf die Gestaltung des eigenen Parts, auf das Verfolgen der
Mitspieler*innen und auf den Orchesterklang.
\autocite[56]{doerne:umfassend_musizieren} Es wird deutlich, dass die
Aufmerksamkeit also auf sehr viele verschiedene Felder gelenkt, sich in
ihren Schwerpunkten immer wieder verschiebt. Aus meiner Sicht
ist dies der faszinierende und aufregende Faktor, der das Orchesterspiel
ausmacht. 

Für Insrumentalist*innen aller Art und jeden Niveaus gibt es Schulorchester,
Kinderorchester, Jugendorchester, Studentenorchester und Laienorchester. Oft
sind diese jedoch in größeren Städten und dafür in ländlicheren Gegenden nicht
ganz so weit verbreitet. Außerdem ist für ein Sinfonieorchester vorausgesetzt,
dass man ein Instrument erlernt hat, welches in der typischen Bestzung eines
Orchesters vorgesehen ist. 



\subsection{Musikverein}

Die Musikvereine sind wiederum vermehrt in ländlicheren Regionen zu Hause. Es
gibt bei ihnen drei verschiedene Instrumental-Kategorien. Zum einen die
Akkordeonvereine, die Blasmusikvereine, sowieo die Fasnachtsgruppierungen. In
den Akkordeonvereinen spielen zwar alle erst einmal das selbe Instrument, aber
doch in ihren unterschiedlichen Stimmen und auch leicht unterschiedlich in der
Funktionsweise, da die Bässe beispielweise verstärkt sind. Manchmal gibt es noch
Trompeter und Schlagzeuger dazu, aber das ist eher die Ausnahme. Bei den
Fasnachtsgruppen, kommen die Angehörigen einer Zunft zusammen und musizieren
gemeinsam. Die Hauptinstrumente hierbei sind Trommeln, Piccolo-Flöten und
Blechblasinstrumente. Sie arbeiten meist nur einmal im Jahr und auf den
Fasnachts-Anlass hin. Im Blasmusikverein sind ähnlich wie im Orchester
verschiedene Instrumentengruppen vertreten aus der sich heraus eine klare
Hierarchie innerhalb der Gruppen bildet. Außerdem gibt es auch eine/n
Dirigent*in, welcher die Proben anleitet. Im Vergleich zum Orchester sind es
aber nur Blasinstrumente, die gespielt werden und ein Schlagzeug zuzüglich. Das
Schöne am Musikverein ist die heterogene Zusammensetzung der Mispieler*innen.
Hier kommen Jung und Alt zusammen, Anfänger und Fortgeschrittene musizieren
gemeinsam. Außer den Landespolizeiorchestern oder beim Militär, gibt es bei
Blasmusikvereinen fast nur Laiengruppierungen. Dadurch, dass also viele
Musiker*innen in ihrer Freizeit im Verein spielen, gehört neben dem eigentlichen
Musizieren auch das Vereinsleben dazu, welches ich von allen hier aufgezählten
Ensembles als das am stärksten ausgeprägteste einordnen würde. Das soziale
Miteinander entsteht somit durch die Musik, aber wird auch darüber hinaus
bewusst gepflegt und gestärkt. Da die Blasmusikvereine eher in ländlicheren
Gegenden verortet werden, spielen sie dort häufig zu besonderen Anlässen im Ort
oder aber auch im Zusammenhang mit Fasnacht zu bestimmten Umzügen. Die Konzerte
sind demnach mit einem festlichen Charakter verbunden, aber trotz allem mit
einem adequaten Grad in der Erwartung an die Leistung der Spieler*innen. Oft
finden Blasmusikvereins-Konzerte auch in einem lockeren Rahmen statt, wo das
Publikum parallel zum Konzert Essen und Trinken bekommt. Das lockert die
Athmosphäre und wird dadurch für das Publikum und auch für die Musizierenden zu
einem gemeinschaftlichen, geselligen Zusammenkommen, welches von Musik begleitet
ist. Wenn man in die Historie zurückblickt, gibt es viele Hauskonzerte und auch
Opernvorstellungen, während denen das Publikum essen, trinken und laut sein
durfte. In diesen Momenten nimmt die Musik eine andere Rolle ein, als in den heutigen
Konzertsälen, in der das Publikum die ernste Musik mit klaren Applausregeln zu
honorieren und ansonsten still zu sein hat. Außerdem ist sie auf diese Art und
Weise sehr natürlich in den Alltag mit eingebunden, während ein
Orchesterkonzertbesuch direkt ein spezieller besonderer Anlass ist. Auch die
Literatur ist in den Musikvereinen näher an den Bedürfnissen der breiten
Gesellschaft, da oft Filmmusik oder andere typische, bekannte Gassenhauer
gespielt werden. Durch die Athmosphäre der Musik wird das Publikum dadurch auch
anders und womöglich besser erreicht, als bspw. bei ernster klassischer Musik. 



\subsection{Chor}

Wie beim Orchester hat ein Chor auch einen/eine Leiter*in, der/die
die Aufwärmübungen und die Probenarbeit mit dem Ensemble anleitet. Parameter,
wie sich in den Klang seiner Gruppe einzuordnen sind auch ähnlich zu anderen
Ensembles, sowie das gemeinsame gleichzeitige und möglichst synchrone Agieren.
Technisch kann man sich jedoch nicht ganz soviel abschauen, da das Singen und
wie es funktioniert nicht so offensichtlich ist, wie das spielen eines
Instrumentes, wo jede Haltung imitiert werden kann. Die Körperhaltung, der Stand
an sich kann beim Singen zwar imitiert werden aber darüber hinaus gilt es seine
Stimmbänder zum klingen zu bringen, was eine sehr individuelle und wenig
sichtbare Angelegenheit darstellt. Musikalisch kann man hingegen viel dazulernen
und gerade die technische Einschränkung die Laien oftmals mit ihren Instrumenten
haben, sind beim Singen etwas weniger eine Barriere. Das ist sicher auch ein
Grund dafür, warum es sehr viele Laienchöre gibt, egal ob in einer größeren
Stadt oder auf dem Dorf. Die Voraussetzungen um in einem Chor mitzusingen sind
oft sehr offen und man muss nie aktiven Gesangsunterricht gemacht haben, um in
einem Laienchor mitwirken zu dürfen. Damit ist das Chor-Ensembles das, was den
einfachsten Zugang zur Musik ermöglicht. 



\subsection{Band}

Mit einer Band assoziiert man direkt verschiedene Genres außerhalb der
klassischen Szene und ein soziales Miteinander, welches im Vergleich zu vielen
anderen Ensemblegruppen nicht unbedingt von einer einzigen Person angeleitet
wird. Im Zusammenhang mit dieser Ensemblearbeit verbinde ich auch das kreative
Arbeiten, die Improvisation und das forschen und ausprobieren, wie man seine
eigenen Stücke mit der Gruppe schreiben, komponieren und umsetzten kann. Da eine
Band aber meist auch nur aus fünf bis zehn Personen besteht, ist diese
Interaktion und der Austausch auch möglich, ohne dass es zu komplex wird. Bei
einer Band sind die Spieler*innen meistens mit verschiedenen Instrumenten am
Musizieren, was ein wichtiger Unterschied zu Orchestern, Chören und
Blasmusikvereinen ist, wo wir zumindest immer eine Gruppering von gleichen
Instrumenten haben. Dadurch, dass jede/r für sein/ihr Instrument verantwortlich
ist und sich die Mitspieler*innen nur mit dem eigenen Instrument auskennen, kann
man sich nur über die Sprache der Musik gegenseitig weiterhelfen. Nicht aber
über konkrete technische Anweisungen. Diese Hürden muss man entweder selbst
überwinden, oder eine Lehrkraft um Hilfe bitten, wenn die Musik als Sprache
selbst nicht mehr weiterhilft. Aber es steht auch viel mehr die Musik selbst im
Vordergrund und es ist zweitrangig, wie man zu seinem Ziel gelangt. Durch die
kreative Herangehensweise müssen die Spieler*innen außerdem eine klare
Vorstellung entwickeln, was sie wie haben möchten und diese wiederum mit ihrem
Mitspielern kommunizieren. Oft haben Musizierende diese klare Vorstellung nicht
von Anfang an, aber es zeichnet Bandproben aus, dass man gemeinsam auf die Suche
danach geht. Dieser Ideenaustausch ist sehr wertvoll und besonders zum
Einzelunterricht eine wunderbare Ergänzung.



\subsection{Sonstige Formationen}

Denkt man an die Instrumente, die auf den ersten Blick nicht für Ensembles bzw.
Gruppenunterrichte geeignet sind, wie zum Beispiel das Klavier oder die Gitarre,
so kommen mir trotzdem und zum Glück doch einige Gruppierungen aus meinem Umfeld
in den Sinn, die sich über die gängigen Formationen nicht abschrecken lassen und
neugierig genug sind, sich ein Ensemble zu suchen. Mit dem Klavier gibt es im
klassischen Sektor wenigstens noch die sparte der Klavier-Kammermusik, von
Klaviertrio bis hin zum Klavierquintett, in anderen Genres ist es aus Bands
nicht wegzudenken. Darüber hinaus wird es eher wieder schwieriger und auch bei
den jungen Klavierschüler*innen, ist das Ensemblespiel leider ein noch viel zu
sehr vernachlässigter Baustein. Wenn man bedenkt welche ganzen sozialen und auch
musikalischen Kompetenzen im Ensemblspiel geschult und gelernt werden, so muss
ich sagen, dass es in Zukunft unbedingt auch für Klavierschüler*innen
Möglichkeiten und Räume zum gemeinsamen Musizieren geschaffen werden müssen!
Betrachten wir die Sparte der neuen Musik, so stelle ich fest, dass gerade hier
immer wieder untypische, neue und besondere Besetzungen von Fagott, Tuba, über Klavier,
Schlagzeug und beispielsweise Violine immer gängiger werden und besonders beliebt
sind. Allerdings sind die Anforderungen an das Niveau in der Neuen Musik meist
leider so hoch, dass sie für Instrumentale Anfänger*innen oft nicht spielbar
sind. Beim Klavier kommt außerdem noch hinzu, dass man aufgrund der Größe des
Instrumentes leider kaum mehr als zwei Klaviere bzw. Flügel in einem Raum stehen
hat. In der Familie der Keyboards sieht es anders aus, da man sie vielseitig
und auch parallell zueinander einsetzen kann.

Vergleichen wir das Klavier beispielsweise mit der Gitarre, die es in der
klassischen Musikszene als Ensembleinstrument auch nicht leicht hat, so kann man
wenigstens mit mehreren Schülern*innen ein Gitarrenensemble zusammenstellen und
gemeinsam musizieren. Dass die Gitarrist*innen aber auch Pianist*innen in den
anderen Genres ihren Platz meist in einer Band haben, vor allem auch im
populären Musikbereich, ist nicht außer Acht zu lassen. Auch in anderen Kulturen
spielt die Gitarre eine sehr ausgeprägte Rolle.

Das Schlagzeug ist ein weiteres Instrument, das es im Ensemble nicht immer
leicht hat. Das Schlagzeug ist ein weiteres Instrument, das zwar in einer
Vielzahl von Ensembleformationen aufzufinden ist, jedoch seltenst in einem
Schlagzeugensemble. Immerhin hat jedes Orchester, jede Band und jeder
Musikverein einen Schlagzeuger. Aber mit seinen eigenen Schlagzeugkolleg*innen
zusammen zu spielen, ist außerhalb des Musikstudiums nicht unbedingt üblich. 

Zusammenfassend kann man sagen, dass die Instrumente in den verschiedenen Genres
andere Funktionen und Einsatzmöglichkeiten haben. Deshalb sollten
Instrumentalschüler*innen an verschiedene Genres herangeführt werden um einen
möglichst breiten Horizont zu erleben und durch diesen die Möglichkeiten sehen,
die ihr Instrument auch im Zusammenhang von Ensembles bietet. 







	\addchap{Fazit}
Ein Plädoyer für das Gemeinsame Musizieren

Wenn wir nun überlegen was wir aus den bestehenden Systemen mitnehmen können, so
komme ich zu dem Ergebnis, dass: 

Dies führt am Ende dazu, dass eine ideale Umgebung für das gemeinsame Musizieren
so geschaffen werden kann, indem man ...

Aufgabe der Musikschule hat der Vdm sehr richtig erkannt, indem er sagt: "(...)
wichtig ist ein ständiges Zusammenwirken von Einzelspiel und Zusammenspiel
einerseits, von praktischem Musizieren und Verstehen des Gespielten
andererseits." (S. 22 die kunst zu unterrichten)

Meiner Meinung nach sollte der Schüler genau dieses Wechselspiel zwischen
Einzelunterricht und gemeinsamen Musizieren in seiner musikalischen Ausbildung
erleben können. Wie wir gelernt haben, fördert es nicht nur die Motivation,
sondern ist auch zur Entwicklung des Einzelnen in seiner Persönlichkeit von
gr0ßer Bedeutung. 


- Ausbildung nicht gut genug S. 119/ S. 94
- EMP hilft und bietet Potential/Techniken

- Schwierigkeit der Organisation

Lösungsansatz:
- Einzelstunden der Schüler kürzen und dafür eine gemeinsame Gruppenstunde im
Anschluss planen (S. 33, die kunst zu unterrichten)
--> Unterrichtszeit für den Lehrer bleibt gleich, aber gerade bei größeren
Gruppen wird die UNterrichtszeit für einen Schüler deutlich länger!
    alternativ kann man auch wöchentlich routieren oder solchen


4 Sozialformen im Instrumentalunterricht S. 34-35    (die Kunst zu unterrichten)

Die wissenschaftliche Erkenntnis als solche reicht nicht aus. Neues Wissen
müssen sich die Lehrenden individuell aneignen und in praktische pädagogische
Kompetenz umwandeln." (S. 10 die Kunst zu unterrichten)

Anselm Ernst S. 61 "Unter dem Motto "von Anfang an" sollten alle Schülerinnen und
Schüler in ein Ensemble aufgenommen werden, ja sogar darauf verpflichtet werde."

- Seite 60-62 erklärt warum Ensemble-Spiel so unumgänglich wichitg ist

	




	\clearpage
	\addcontentsline{toc}{chapter}{Literatur}
	\printbibliography

	% \clearpage
	% \addcontentsline{toc}{chapter}{Stichwortverzeichnis}
	% \printindex

	% \clearpage
	% \addcontentsline{toc}{chapter}{Verzeichnis der Notenbeispiele}
	% \listoffigures

	\section*{Erklärung}

\pagenumbering{gobble}

Hiermit erkläre ich, dass ich diese Arbeit selbständig verfasst habe, keine
anderen als die angegebenen Quellen und Hilfsmittel verwendet habe und alle
Stellen, die wörtlich oder sinngemäß aus veröffentlichten Schriften entnommen
wurden, als solche kenntlich gemacht habe. Darüber hinaus erkläre ich, dass
diese Arbeit nicht, auch nicht auszugsweise, bereits für eine andere Prüfung
angefertigt wurde.


\end{document}
